% ═══════════════════════════════════════════════════════════════════════════════
%  CAPÍTULO — Adapter
% ═══════════════════════════════════════════════════════════════════════════════

\chapter{Adapter}

\patroninfo{Estructural}{139}

% ─── Situación Problema ──────────────────────────────────────────────────────
\section{Situación Problema}

Airbnb integra pasarelas de pago de terceros —Stripe, PayPal y
Mercado Pago— para procesar los cobros a los huéspedes. Cada pasarela
expone su propia interfaz de programación: Stripe utiliza un método
\texttt{charge()}, PayPal emplea \texttt{executePayment()} y Mercado
Pago trabaja con \texttt{createPaymentIntent()}, todos con firmas,
parámetros y formatos de respuesta distintos.

El servicio de pagos interno de Airbnb no puede —ni debería— adaptarse
a la interfaz de cada proveedor. Si en el futuro se incorpora una nueva
pasarela, habría que modificar el servicio central para entender su
vocabulario, acoplando el núcleo del negocio a detalles de implementación
externos y fragilizando el sistema.

% ─── Solución ────────────────────────────────────────────────────────────────
\section{Solución}

El patrón Adapter introduce una interfaz unificada \texttt{PasarelaPago}
con un único método \texttt{procesarPago()}. Para cada proveedor externo
se crea un adaptador concreto —\texttt{AdaptadorStripe},
\texttt{AdaptadorPayPal}, \texttt{AdaptadorMercadoPago}— que implementa
esa interfaz y traduce internamente la llamada al método nativo de cada
SDK.

El servicio de pagos de Airbnb trabaja exclusivamente contra
\texttt{PasarelaPago} y desconoce por completo las particularidades de
cada proveedor. Añadir una nueva pasarela equivale únicamente a crear
un nuevo adaptador, sin tocar el código de negocio. Esto también
facilita las pruebas unitarias, pues basta con sustituir el adaptador
real por uno simulado.

% ─── Diagrama de clases ─────────────────────────────────────────────────────
\begin{figure}[H]
    \centering
    % \includegraphics[width=\textwidth]{imagenes/abstract_factory_diagrama.png}
    \caption{Diagrama de clases del patrón Abstract Factory}
\end{figure}


% ─── Roles de las Clases ────────────────────────────────────────────────────
\section{Roles de las Clases}

% Explicar la responsabilidad de cada clase/interfaz

\begin{itemize}
    \item \textbf{AbstractFactory:}
    \item \textbf{ConcreteFactory:}
    \item \textbf{AbstractProduct:}
    \item \textbf{ConcreteProduct:}
    \item \textbf{Client:}
\end{itemize}


% ─── Main ───────────────────────────────────────────────────────────────────
\section{Main}

% Explicar qué realiza el programa principal

\begin{lstlisting}[language=Java, caption=NombreClase]
 
\end{lstlisting}


% ─── Salida del Main ────────────────────────────────────────────────────────
\section{Salida del Main}

\begin{lstlisting}[language=Java, caption=NombreClase]
 
\end{lstlisting}


% ─── Clases ─────────────────────────────────────────────────────────────────
\section{Clases}

% ─── Clase 1 ────────────────────────────────────────────────────────────────
\subsection{NombreClase}

\begin{lstlisting}[language=Java, caption=NombreClase]
 
\end{lstlisting}


% ─── Clase 2 ────────────────────────────────────────────────────────────────
\subsection{NombreClase}

\begin{lstlisting}[language=Java, caption=NombreClase]
 
\end{lstlisting}


% ─── Clase 3 ────────────────────────────────────────────────────────────────
\subsection{NombreClase}

\begin{lstlisting}[language=Java, caption=NombreClase]
 
\end{lstlisting}


% ─── Conclusiones ───────────────────────────────────────────────────────────
\section{Conclusiones}

% Explicar:
% - Ventajas del patrón
% - Cuándo usarlo
% - Beneficios obtenidos
% - Posibles desventajas
% - Relación con principios SOLID